\chapter{Introduction}
\label{ch:introduction}

\section{The Puzzle of Post-\textit{Dobbs} Diffusion}

On June 24, 2022, the Supreme Court issued its decision in \textit{Dobbs v. Jackson Women's Health Organization}, overruling \textcite{RoeWade1971} and \textcite{PlannedParenthoodSoutheastern1992} and returning the regulation of abortion to the states. The decision did not merely alter constitutional doctrine. It collapsed, in a single opinion, the federal preemption regime that had structured state-level reproductive policy for nearly half a century. The states became, overnight, the exclusive jurisdictional venue for a contested policy domain in which both regulatory expansion and codification of protections had been institutionally constrained for two generations.

The legislative response was rapid, substantively bilateral, and---I will argue---structurally novel. By the close of the 2023 legislative session, forty-six states had enacted at least one statute responsive to the new constitutional environment, producing a corpus of 188 bills that ranged from six-week gestational limits and trigger statutes to shield laws, codified clinic protections, and procedural expansions of access. The textual record of that single legislative year is the empirical object I analyze.

The accumulated literature on policy diffusion did not anticipate the characteristic features of this record. The geographic-learning tradition, descending from \textcite{walkerDiffusionInnovationsAmerican1969}, predicts that innovation spreads through regional clusters as adopting states observe and adapt to the experience of nearby pioneers. The event-history framework of \textcite{berryStateLotteryAdoptions1990} specifies internal determinants and regional influence as the principal predictors of adoption timing. The copy-and-paste tradition formalized by \textcite{jansaCopyPasteLawmaking2019} predicts that interest-group drafting infrastructure produces detectable textual signatures, and---given the institutional prominence of national restrictive-coalition organizations in pre-\textit{Dobbs} abortion policy---should produce its strongest signature on the restrictive side. None of these predictions accurately describe the 2023 record. Geographic contiguity is statistically null as a predictor of cross-state textual similarity. Coalition coordination is detectable, but it is concentrated on the protective rather than the restrictive side. And the diffusion timeline is too compressed for the observational learning the geographic tradition presupposes: bills enacted in March cannot be the product of states having watched the consequences of bills enacted in February.

The empirical puzzle, then, concerns diffusion regimes. When federal preemption abruptly collapses in a highly salient, organizationally mature policy domain, what mechanism replaces the slow regional learning the field has theorized? The answer I develop is that the mechanism is not learning at all. It is the activation of pre-positioned templates through coalition-aligned organizational networks---a regime that the diffusion literature has not previously needed to specify, and one whose signatures are visible in the textual record of the first post-shock legislative cycle.

\section{The Argument}

I advance four interlocking claims. The first is empirical: the 2023 post-\textit{Dobbs} abortion corpus exhibits a characteristic textual structure---moderate cross-state borrowing rather than either wholesale copying or independent drafting, with two-thirds of bills clustered between 0.70 and 0.90 on the cross-state novelty measure---and that structure is coalition-asymmetric. Protective bills are significantly less novel than the neutral baseline ($\beta = -0.076$, $p = 0.022$); restrictive bills are not statistically distinguishable from it ($\beta = -0.017$, $p = 0.482$). At the dyadic level, cross-state textual similarity is predicted by shared organizational infrastructure ($\beta = +0.039$, $p < 0.001$, measured as the Jaccard index of PAC donor organizations) and by shared protective coalition membership ($\beta = +0.037$, $p < 0.001$), not by geographic contiguity ($\beta = +0.003$, $p = 0.734$). These findings invert the pre-\textit{Dobbs} expectation that restrictive-coalition templating would dominate the observable record.

The second claim is theoretical. The thesis develops a framework---\emph{constitutional vacuum diffusion}---that specifies the structural conditions under which the empirical signature documented here should be expected: a judicially imposed rather than discursively negotiated policy window, a bilateral rather than unilateral opportunity structure, a compressed temporal horizon, and a pre-existing advocacy infrastructure that positioned model language before the precipitating shock. Within that framework, the thesis identifies \emph{reactive template mobilization} as the operating mechanism: pre-built statutory language, refined over years by organizations on both sides of the issue, is activated through affiliate channels in the months immediately following the shock. The mechanism is distinct from learning-based diffusion (which presupposes time to observe consequences), from conventional crisis adoption (which presupposes templates produced in response to the precipitating event), and from the Shipan and Volden four-mechanism typology (whose categories the post-shock hybrid does not cleanly match).

The third claim concerns coalition architecture. The asymmetry between restrictive and protective coordination is not, on the evidence assembled here, an accident of mobilization timing. It reflects two distinct organizational forms: a vertical architecture on the restrictive side, in which template development concentrates within a small number of national nodes (Americans United for Life, Susan B. Anthony Pro-Life America, the Alliance Defending Freedom) that distribute language outward to allied state legislators; and a horizontal architecture on the protective side, in which state-level affiliates of the principal advocacy organizations (Planned Parenthood, the ACLU, the Center for Reproductive Rights, Reproductive Freedom for All (RFA)) coordinate laterally on a shared codification agenda. These architectures produce different textual signatures, and the 2023 record is consistent with both operating simultaneously.

The fourth claim is methodological. The dissociation between financial transfer and textual coordination---reported across the ten highest-similarity bill pairs and confirmed by systematic analysis of the DIME contribution record---provides empirical support for the \textcite{hallLobbyingLegislativeSubsidy2006a} legislative-subsidy framework in a policy domain where the financial-influence intuition has been particularly prominent in public commentary. The organizations whose Jaccard overlap predicts convergence are not purchasing legislative behavior. They are supplying informational infrastructure: model language, legal memoranda, drafting assistance, and expert testimony. Contribution-only measures of interest-group influence systematically misspecify the channel through which coordination operates in this domain.

\section{Scope, Approach, and Limits}

This thesis was designed to answer a narrow question with precision: what textual and structural patterns characterize cross-state diffusion in the first post-\textit{Dobbs} legislative cycle, and which institutional and coalitional features predict them? The empirical scope is correspondingly constrained. The analysis covers a single legislative year (2023), a single policy domain (abortion regulation), and a single national context (the United States). It analyzes enacted legislation rather than the full population of introduced bills. It measures textual similarity through TF--IDF cosine distance over $n$-gram representations rather than through semantic embeddings. And it operationalizes organizational infrastructure through PAC donor co-presence rather than through richer measures of organizational depth that would require lobby-disclosure and IRS data, the thesis does not incorporate.

These choices were deliberate prioritizations rather than concessions, a framing developed at length in Section~\ref{sec:designed-scope}. The single-year scope isolates the initial legislative response to the constitutional shock before feedback effects from judicial challenges and electoral consequences could contaminate the diffusion signal. The enacted-bill scope captures the completed output of the legislative process and avoids the interpretive ambiguity of bills whose failure may reflect strategic withdrawal rather than substantive rejection. The TF--IDF representation provides transparent, reproducible textual measurement consistent with the broader text-as-data literature on legislative similarity \parencite{wilkersonTracingFlowPolicy2015,linderTextPolicyMeasuring2020}. Each choice carries inferential costs, and Chapter~\ref{ch:limitations} specifies them explicitly along with the future research agenda---qualitative inquiry into drafting processes, archival enrichment of organizational measures, longitudinal corpus extension, cross-domain replication, and semantic methodological refinement---that would address them.

The argument's implication structure extends well beyond the narrow empirical question. If constitutional vacuum diffusion is a generalizable regime rather than a description of one policy area's idiosyncratic post-shock dynamics, the framework should travel to other domains in which federal doctrinal infrastructure has collapsed or may collapse---firearms regulation following \textit{New York State Rifle \& Pistol Association v. Bruen}, voting rights following \textit{Shelby County v. Holder}, hypothetical reversals of other long-standing constitutional doctrines. Whether the framework does travel is an empirical question that the present thesis raises but does not answer. The contribution it makes is to specify the framework precisely enough that subsequent research can evaluate it.

\section{Plan of the Thesis}

The chapters that follow proceed from the literatures the thesis draws on through the empirical analysis to the theoretical apparatus and its limits.

Chapter~\ref{ch:theory} reviews the diffusion literature that the thesis engages and revises: the geographic tradition, the event-history framework, the copy-and-paste research program, the state-capture literature on organizational influence, and the four-mechanism typology of Shipan and Volden. It identifies the specific predictions each tradition generates for the post-\textit{Dobbs} environment and locates the analytical gap that constitutional vacuum diffusion is positioned to fill.

Chapter~\ref{ch:methodology} specifies the research design. It documents corpus construction from LegiScan records, the TF--IDF representation and cosine-similarity computation, the construction of the cross-state novelty measure, the bill-level regression specification, and the dyadic regression specification incorporating the PAC Jaccard organizational measure derived from the DIME database. It justifies the analytical choices the design makes and previews the limitations the design imposes.

Chapter~\ref{ch:findings} reports the empirical findings. It documents the distributional properties of the novelty measure, the bill-level regression results including the partial inversion of \textcite{squireMeasuringStateLegislative2007} on professionalism, and the dyadic regression results that establish organizational co-presence and coalition alignment as the structural predictors of cross-state textual similarity. It presents the case-level evidence of the Nevada SB 370--Washington HB 1155 near-replication and the dissociation between financial transfer and textual coordination across the highest-similarity bill pairs.

Chapter~\ref{ch:discussion} develops the theoretical apparatus. It introduces constitutional vacuum diffusion as a regime characterized by specific structural conditions, specifies reactive template mobilization as the operating mechanism within that regime, develops the vertical--horizontal coalition architecture distinction between restrictive and protective coordination, reinterprets the professionalism finding as a measure of capacity for adoption rather than capacity for innovation, and reconstructs the exporter--importer--reinforcer typology to accommodate the 2023 record's coalition-asymmetric structure.

Chapter~\ref{ch:limitations} specifies the limits of the design. It identifies the inferential gap between the textual signatures the analysis detects, and the mechanism-level claims the framework advances, articulates the rival mechanisms the design cannot eliminate, and develops the future research agenda---qualitative drafting-process inquiry, archival enrichment, longitudinal extension, cross-domain replication, and semantic methodological refinement---that subsequent investigation should pursue.

Chapter~7 consolidates the contributions, locates the thesis within the trajectory of diffusion scholarship, and closes with the interpretive question the 2023 record poses most urgently for the field: whether the patterns documented here represent a transitional anomaly of the immediate post-\textit{Dobbs} mobilization period or the first evidence of a structural transformation in how American states make policy after constitutional shocks. That question cannot be resolved from a single legislative year. The analytical infrastructure assembled here is designed to support its resolution by the research that follows.